\clearpage
\onecolumn
\raggedbottom
\section{Agent API and Control Surface Reference}
\label{app:agent-api}

\newcommand{\apicmd}[1]{\texttt{\footnotesize #1}}
\newcommand{\apicode}[1]{\texttt{\footnotesize #1}}
\newcommand{\apitabletitle}[1]{\par\medskip\begin{center}\small\textsc{#1}\end{center}\vspace{-0.7em}}
\setlength{\tabcolsep}{3pt}
\renewcommand{\arraystretch}{1.12}

This appendix is the agent-facing contract for the V3.01 repository. It is grounded in the active \apicode{jeryu} binary, the Unix-socket capability API, the native MCP server, the engine HTTP routes, the GitLab Runner custom executor protocol, the Git server admission hook, and the canonical action registry. The MCP surface is not a parallel authority system: it is a standards-compatible transport for the same policy, evidence, grant, and merge/release checks exposed through the existing capability model.

Git compatibility is handled through \apicode{jeryu git <args>} passthrough plus the native \apicode{jeryu save}, \apicode{jeryu sync}, and \apicode{jeryu undo} wrappers. The removed \apicode{ship} command and the old repo-mirror entrypoints no longer own a separate jeryu-git control plane.

\subsection{Surface Map}

\apitabletitle{Agent-Relevant API Surfaces}
\begin{center}
\small
\begin{tabular}{@{}p{0.15\linewidth}p{0.13\linewidth}p{0.20\linewidth}p{0.16\linewidth}p{0.24\linewidth}@{}}
\toprule
\textbf{Surface} & \textbf{Transport} & \textbf{Primary Consumer} & \textbf{Auth/Gate} & \textbf{Use}\\
\midrule
CLI & process invocation & humans, local agents, scripts & local process authority, GitLab PAT where needed & lifecycle, tests, release, pools, evidence, host operations\\
TUI & terminal UI & human supervisors, screen-reading agents & local process authority & operational state, actions, logs, proof stack\\
Capability API & Unix domain socket JSON & supervised local agents & protocol envelope, nonce, optional grant proof & typed agent intents and action discovery\\
MCP server & stdio JSON-RPC; loopback Streamable HTTP & MCP-compatible local agents & MCP lifecycle plus capability policy & standard tool discovery and calls over the same grant/evidence gates\\
HTTP engine & HTTP & GitLab webhooks, local health checks & webhook token for \apicode{/hooks} & reconciliation, cache summary, webhooks\\
Server hook & stdin/stdout process hook & Git server / GitLab hook & admission ledger, optional enforce mode & pre-receive branch/ref policy\\
Custom executor & GitLab Runner protocol & runner managers & runner invocation contract & sandbox, cache, witness, failure capsule capture\\
Action registry & static JSON contract & CLI, TUI, capability clients & canonical risk/grant metadata & discover actions and required grants\\
\bottomrule
\end{tabular}
\end{center}

\begin{lstlisting}[style=api,language=bash,caption={Agent discovery and status entrypoints}]
jeryu action list --json
jeryu status
jeryu next --project-id 2 --ref-name main
jeryu explain-blocker job 14445
jeryu capability serve /tmp/jeryu-capability.sock
jeryu mcp serve
jeryu mcp serve-http
jeryu mcp tools --json
\end{lstlisting}

\clearpage
\subsection{CLI Command Catalog}

All command names are kebab-case at the shell. Defaults below match the current Clap definitions.

\apitabletitle{Lifecycle, TUI, Status, and Global Action Commands}
\begin{center}
\small
\begin{tabular}{@{}p{0.28\linewidth}p{0.13\linewidth}p{0.53\linewidth}@{}}
\toprule
\textbf{Command} & \textbf{Mode} & \textbf{Purpose}\\
\midrule
\apicmd{jeryu init} & mutating & Bootstrap GitLab, state DB, runners, secrets, compose files, and smoke setup.\\
\apicmd{jeryu bootstrap} & mutating & Hidden alias for \apicode{init}.\\
\apicmd{jeryu serve} & daemon & Start Docker/GitLab orchestration, hook install, cache supervisor, runner warmup, engine, and background workers.\\
\apicmd{jeryu down} & mutating & Drain runner managers and stop the GitLab stack.\\
\apicmd{jeryu status} & read & Show GitLab, Vault, pools, managed containers, jobs, release, and secret state.\\
\apicmd{jeryu git <args>} & mixed & Pass through to the system Git binary and preserve the hook-driven jeryu-git path.\\
\apicmd{jeryu save|sync|undo} & mixed & Convenience wrappers for add/commit, rebase+push, and soft reset.\\
\apicmd{jeryu tui [--once]} & UI/read & Launch the Ratatui dashboard or render one smoke frame.\\
\apicmd{jeryu tui --capture --tab <tab> --output <png>} & read/artifact & Deterministic Ratatui PNG capture. Tabs: mission, release, jobs, agents, tests, pools, cache, evidence, secrets.\\
\apicmd{jeryu tui --screenshot --tab <tab> --screenshot-hold-ms <ms>} & read/session & Hold a real terminal render for external screenshot tooling.\\
\apicmd{jeryu progress --project-id 2 --ref-name main [--json]} & read & Show lane-aware CI and release progress.\\
\apicmd{jeryu next --project-id 2 --ref-name main} & read & Print the highest-priority next action for a branch.\\
\apicmd{jeryu explain-blocker <job|release|merge> <id>} & read & Explain why an entity is blocked using DB/GitLab evidence.\\
\apicmd{jeryu action list [--json]} & read & Emit the canonical action registry for humans or agents.\\
\bottomrule
\end{tabular}
\end{center}

\clearpage
\apitabletitle{Pool, Job, Pipeline, Cache, and Host Commands}
\begin{center}
\small
\begin{tabular}{@{}p{0.32\linewidth}p{0.13\linewidth}p{0.51\linewidth}@{}}
\toprule
\textbf{Command} & \textbf{Mode} & \textbf{Purpose}\\
\midrule
\apicmd{jeryu pool list} & read & List pools, manager counts, paused state, executor, and runner id.\\
\apicmd{jeryu pool scale <name> <count>} & mutating & Scale a pool to an exact manager count.\\
\apicmd{jeryu pool pause|resume <name>} & mutating & Pause or resume runner intake.\\
\apicmd{jeryu pool drain <name>} & mutating & Pause, wait, and stop managers.\\
\apicmd{jeryu pool delete <name>} & mutating & Drain, delete DB record, and unregister runner.\\
\apicmd{jeryu pool rotate-token <name>} & mutating & Reset GitLab runner token and update manager config.\\
\apicmd{jeryu job list <project> [--status running,pending]} & read & List GitLab jobs by scope.\\
\apicmd{jeryu job trace <project> <job>} & read & Print GitLab job trace.\\
\apicmd{jeryu job play|cancel|retry <project> <job>} & mutating & Start manual jobs, cancel jobs, or retry failed jobs.\\
\apicmd{jeryu job explain <project> <job>} & read & Print latest structured failure capsule and retry advice.\\
\apicmd{jeryu job clear} & mutating & Clear local job and pipeline history.\\
\apicmd{jeryu pipeline explain|doctor|jobs <pipeline> [--project-id 2] [--json]} & read & Inspect release eligibility, health, and job timing.\\
\apicmd{jeryu pipeline jobs <pipeline> --ingest} & mutating & Store current job timing rows.\\
\apicmd{jeryu pipeline ingest|cancel <pipeline> [--project-id 2]} & mutating & Persist timings or cancel a pipeline.\\
\apicmd{jeryu pipeline bottlenecks [--limit 25] [--json]} & read & Report historical slow jobs from local timing ledger.\\
\apicmd{jeryu cache enable} & mutating & Configure Docker for SmartCache registry mirror.\\
\apicmd{jeryu cache doctor|status [--json]} & read & Check proxy/registry/cache state and metrics.\\
\apicmd{jeryu cache gc [--dry-run] [--json] [--older-than <age>] [--max-cache-gb <gb>]} & mutating & Garbage-collect manager caches and cache state.\\
\apicmd{jeryu logs <manager-id> [-n 50]} & read & Tail manager container logs.\\
\apicmd{jeryu host storage-audit|doctor [--json]} & read & Audit storage and host health.\\
\apicmd{jeryu host reclaim --mode <mode> (--plan|--apply)} & mutating & Run planned or applied reclaim.\\
\bottomrule
\end{tabular}
\end{center}

\clearpage
\apitabletitle{Agent, Jeryu, Test/VTI, Release, Secret, and Repo Commands}
\begin{center}
\small
\begin{tabular}{@{}p{0.33\linewidth}p{0.13\linewidth}p{0.50\linewidth}@{}}
\toprule
\textbf{Command} & \textbf{Mode} & \textbf{Purpose}\\
\midrule
\apicmd{jeryu agent spawn <project> --task <text>} & mutating & Create an autonomous task branch/MR/issue workflow.\\
\apicmd{jeryu agent list <project>} & read & List active agent MRs/issues.\\
\apicmd{jeryu agent merge <project> <mr> [--trust-tier trusted]} & mutating & Request merge through the risk gate.\\
\apicmd{jeryu test run --command <cmd> [--force]} & mutating & Run a command through an ephemeral GitLab CI pipeline.\\
\apicmd{jeryu test plan --command <cmd>} & read & Preview inferred image, tags, timeout, and runner class.\\
\apicmd{jeryu test batch --command <cmd>... [--max-parallel 3]} & mutating & Launch several validation pipelines concurrently.\\
\apicmd{jeryu test results|failed|retry <pipeline>} & mixed & Inspect or retry pipeline test jobs.\\
\apicmd{jeryu test impact --base <sha> --head <sha> [--json]} & read & Determine CI jobs and release gates needed for a diff.\\
\apicmd{jeryu test select [emit flags]} & read/artifact & Compute internal VTI plan and optionally emit GitLab YAML, JSON plan, and VTI receipt.\\
\apicmd{jeryu test explain-plan <plan.json>} & read & Explain a serialized VTI plan.\\
\apicmd{jeryu test select-external --workspace <root>} & read/artifact & Use external \apicode{.jeryu/testmap.toml}.\\
\apicmd{jeryu test audit|learn} & read & Compare VTI selections with oracle results and suggest rule changes.\\
\apicmd{jeryu test cache-status --base <sha> --head <sha> [--json]} & read & Show test-command cache state for the current source state.\\
\apicmd{jeryu release status|watch|reconcile} & mixed & Read or reconcile release attempts.\\
\apicmd{jeryu release promote-prod [--version <v>]} & mutating & Trigger approved production promotion after canary gates.\\
\apicmd{jeryu release preflight|doctor [--json]} & read & Check blockers, infrastructure, and release readiness.\\
\apicmd{jeryu secrets <subcommand>} & mixed & Manage Vault-backed release lifecycle: init, status, rotate, finalize, report, recover.\\
\apicmd{jeryu repo render-agent-index [--check]} & artifact/check & Generate or verify machine-readable agent routing index.\\
\apicmd{jeryu repo audit-agent-surface [--json]} & read & Audit routing docs, generated index, and agent surface freshness.\\
\bottomrule
\end{tabular}
\end{center}

\subsection{Exact CLI Grammar}

The following grammar is generated from the active Clap definitions in \apicode{src/cli.rs}. It includes hidden protocol calls because those are part of the agent and runner contract even though they are not intended for direct human operation.

\begin{lstlisting}[style=api,language=bash,caption={Exhaustive jeryu command forms}]
jeryu init
jeryu bootstrap
jeryu serve
jeryu down
jeryu status
jeryu logs <manager-id> [-n <lines>]
jeryu progress [--project-id <id>] [--ref-name <ref>] [--json]
jeryu next [--project-id <id>] [--ref-name <ref>]
jeryu explain-blocker <job|release|merge> <id>

jeryu tui [--once]
jeryu tui --capture [--tab <tab>] [--output <path>] [--width <cols>] [--height <rows>]
jeryu tui --screenshot [--tab <tab>] [--width <cols>] [--height <rows>] [--screenshot-hold-ms <ms>]

jeryu action list [--json]

jeryu pool list
jeryu pool scale <name> <count>
jeryu pool pause <name>
jeryu pool resume <name>
jeryu pool drain <name>
jeryu pool delete <name>
jeryu pool rotate-token <name>

jeryu job list <project-id> [--status <csv>]
jeryu job trace <project-id> <job-id>
jeryu job play <project-id> <job-id>
jeryu job cancel <project-id> <job-id>
jeryu job retry <project-id> <job-id>
jeryu job explain <project-id> <job-id>
jeryu job clear

jeryu pipeline explain [--project-id <id>] <pipeline-id> [--json]
jeryu pipeline doctor [--project-id <id>] <pipeline-id> [--json]
jeryu pipeline jobs [--project-id <id>] <pipeline-id> [--ingest] [--json]
jeryu pipeline ingest [--project-id <id>] <pipeline-id>
jeryu pipeline cancel [--project-id <id>] <pipeline-id>
jeryu pipeline bottlenecks [--project-id <id>] [--ref-name <ref>] [--limit <n>] [--json]

jeryu cache enable
jeryu cache doctor
jeryu cache status [--json]
jeryu cache gc [--dry-run] [--json] [--keep-active-managers[=true|false]] [--older-than <age>] [--max-cache-gb <gb>]

jeryu agent spawn <project-id> --task <text>
jeryu agent list <project-id>
jeryu agent merge <project-id> <mr-iid> [--trust-tier <tier>]

jeryu test run --command <cmd> [--project-id <id>] [--image <image>] [--tags <csv>] [--timeout <sec>] [--force]
jeryu test plan --command <cmd> [--project-id <id>] [--image <image>] [--tags <csv>] [--timeout <sec>]
jeryu test batch --command <cmd>... [--project-id <id>] [--image <image>] [--tags <csv>] [--timeout <sec>] [--max-parallel <n>] [--force]
jeryu test results <pipeline-id> [--project-id <id>]
jeryu test retry <pipeline-id> <job-name> [--project-id <id>]
jeryu test failed <pipeline-id> [--project-id <id>]
jeryu test impact --base <ref> --head <ref> [--repo-root <path>] [--json]
jeryu test select [--base <ref>] [--head <ref>] [--repo-root <path>] [--explain] [--json] [--emit-gitlab <path>] [--emit-plan <path>] [--emit-receipt <path>]
jeryu test explain-plan <plan-path>
jeryu test select-external [--base <ref>] [--head <ref>] --workspace <path> [--explain] [--json] [--emit-gitlab <path>] [--emit-plan <path>] [--emit-skipped <path>]
jeryu test audit --changed <csv> --all-tests <csv> [--failed <csv>] [--sha <sha>] [--json] [--workspace <path>]
jeryu test learn --changed <csv> --all-tests <csv> [--failed <csv>] [--sha <sha>] [--json] [--workspace <path>]
jeryu test cache-status [--base <ref>] [--head <ref>] [--json]

jeryu release status [--project-id <id>] [--ref-name <ref>] [--sha <sha>] [--limit <n>] [--json]
jeryu release watch [--project-id <id>] [--ref-name <ref>] [--sha <sha>] [--limit <n>] [--interval-secs <sec>] [--json]
jeryu release reconcile [--project-id <id>] [--ref-name <ref>] [--json]
jeryu release promote-prod [--project-id <id>] [--ref-name <ref>] [--version <version>]
jeryu release preflight [--ssh-host <host>] [--json]
jeryu release doctor [--version <version>] [--preflight[=true|false]] [--json]

jeryu secrets init
jeryu secrets status [--json]
jeryu secrets rotate [--repo <name>] --version <version> --target <target>
jeryu secrets finalize [--repo <name>] --version <version> --target <target>
jeryu secrets report [--repo <name>] --version <version>
jeryu secrets recover [--repo <name>] --version <version>

jeryu host storage-audit
jeryu host doctor [--json]
jeryu host reclaim --mode <mode> [--plan] [--apply]

jeryu repo render-agent-index [--check]
jeryu repo audit-agent-surface [--json]

jeryu exec config
jeryu exec prepare
jeryu exec run <script-path> <stage>
jeryu exec cleanup
jeryu server-hook pre-receive
jeryu capability serve <socket-path>
\end{lstlisting}

\clearpage
\subsection{Hidden Protocol Entrypoints}

\apitabletitle{Protocol Commands Not Intended for Direct Human Use}
\begin{center}
\small
\begin{tabular}{@{}p{0.29\linewidth}p{0.17\linewidth}p{0.49\linewidth}@{}}
\toprule
\textbf{Command} & \textbf{Caller} & \textbf{Contract}\\
\midrule
\apicmd{jeryu exec config} & GitLab Runner & Return custom executor driver configuration JSON.\\
\apicmd{jeryu exec prepare} & GitLab Runner & Prepare clone/sandbox/honeypot/cache context.\\
\apicmd{jeryu exec run <script> <stage>} & GitLab Runner & Run a build stage with cache brain, witnesses, tripwire handling, and capsule capture.\\
\apicmd{jeryu exec cleanup} & GitLab Runner & Clean executor state after job completion.\\
\apicmd{jeryu server-hook pre-receive} & Git server hook & Read ref updates from stdin and write admission decisions.\\
\apicmd{jeryu capability serve <socket>} & local supervisor & Start Unix-socket JSON capability API.\\
\bottomrule
\end{tabular}
\end{center}

\subsection{Capability Socket Protocol}

The capability API is a local Unix domain socket service. It accepts either a  tagged \apicode{AgentIntent} JSON document or the V3.01 \apicode{AgentActionRequest} envelope. New agents should use the envelope. The server also accepts an optional 4-byte big-endian length prefix before the JSON body. Responses use \apicode{CapabilityResponse}.

\begin{lstlisting}[style=api,language=rust,caption={V3.01 capability envelope}]
struct AgentActionRequest {
  protocol_version: "v3.01",
  request_id: String,
  actor: String,
  nonce: String,
  expires_at: Option<Rfc3339DateTime>,
  project_id: Option<i64>,
  base_ref: Option<String>,
  base_sha: Option<String>,
  idempotency_key: Option<String>,
  budget: Option<ActionBudget>,
  grant: Option<CapabilityGrantProof>,
  intent: AgentIntent
}

struct CapabilityResponse {
  success: bool,
  message: String,
  data: Option<JsonValue>
}
\end{lstlisting}

\begin{lstlisting}[style=api,language=json,caption={Envelope example: request a system snapshot}]
{
  "protocol_version": "v3.01",
  "request_id": "req-20260427-0001",
  "actor": "agent.codex.local",
  "nonce": "7f096ab2-9e0d-4d80-9c4b-42fb56c6d62d",
  "expires_at": "2026-04-27T10:15:00Z",
  "idempotency_key": "snapshot-main-001",
  "intent": { "intent": "GetSystemSnapshot" }
}
\end{lstlisting}

\apitabletitle{Envelope Validation and Replay Rules}
\begin{center}
\small
\begin{tabular}{@{}p{0.22\linewidth}p{0.72\linewidth}@{}}
\toprule
\textbf{Field/Rule} & \textbf{Behavior}\\
\midrule
\apicode{protocol\_version} & Must equal \apicode{v3.01}; other versions are rejected.\\
\apicode{request\_id} and \apicode{actor} & Required for enveloped requests and logged as request context.\\
\apicode{nonce} & Required and cached per actor; duplicate actor/nonce pairs are rejected as replay.\\
\apicode{expires\_at} & Optional RFC3339 timestamp; expired requests are rejected.\\
\apicode{budget} & Optional caller-declared time/branch/pipeline limits.\\
\apicode{grant} & Optional proof payload. Current code validates envelope shape and records grants for successful branch-writing intents; signed grant enforcement remains future hardening.\\
Legacy requests & Accepted for compatibility, tagged as \apicode{-capability-client}.\\
\bottomrule
\end{tabular}
\end{center}

\subsection{Capability Intents}

\apitabletitle{Capability Intents and Effects}
\begin{center}
\small
\begin{tabular}{@{}p{0.18\linewidth}p{0.15\linewidth}p{0.22\linewidth}p{0.41\linewidth}@{}}
\toprule
\textbf{Intent} & \textbf{Risk} & \textbf{Inputs} & \textbf{Result}\\
\midrule
\apicode{FetchCapsule} & read & \apicode{job\_id} & Return latest DB failure capsule for a job, or a structured miss.\\
\apicode{RunTests} & CI execution & \apicode{project\_id}, \apicode{target\_ref}, \apicode{test\_scope} & Create ephemeral branch, write dynamic CI YAML, trigger pipeline, record a 24-hour grant.\\
\apicode{ProposePatch} & git write & project, branch, base, commit message, file modifications, optional MR title & Create branch, commit changes, open MR, record grant bound to returned commit SHA.\\
\apicode{RacePatches} & git write / CI & project, base branch, commit message, hypotheses & Launch multiple hypothesis branches and pipelines; returns branches, pipeline ids, grant ids.\\
\apicode{RequestMerge} & production/merge & project, MR IID, source, target & Evaluate conservative merge gate proof using selector miss and cache taint checks.\\
\apicode{ExplainBlockers} & read & entity type and id & Explain blockers for job, release, or merge-like entity.\\
\apicode{GetSystemSnapshot} & read & none & Return pools, cache metrics, selector misses, recent failures, timestamp.\\
\apicode{ListAllowedActions} & read & none & Return capability-visible registry actions with risk/grant metadata.\\
\apicode{PlanValidation} & read & project, test ids, ref & Validate a proposed test list against recent selector misses.\\
\bottomrule
\end{tabular}
\end{center}

\begin{lstlisting}[style=api,language=json,caption={Legacy intent examples}]
{ "intent": "FetchCapsule", "payload": { "job_id": 14445 } }

{
  "intent": "RunTests",
  "payload": {
    "project_id": 2,
    "target_ref": "main",
    "test_scope": "unit"
  }
}

{
  "intent": "PlanValidation",
  "payload": {
    "project_id": 2,
    "ref_name": "main",
    "test_ids": ["unit:jeryu", "integration:job_tests"]
  }
}
\end{lstlisting}

\begin{lstlisting}[style=api,language=json,caption={ProposePatch payload shape}]
{
  "intent": "ProposePatch",
  "payload": {
    "project_id": 2,
    "branch_name": "agent/fix-cache-proof",
    "base_ref": "main",
    "commit_message": "fix: strengthen cache proof",
    "mr_title": "Agent patch: cache proof hardening",
    "modifications": [
      {
        "file_path": "src/cache_brain.rs",
        "content": "...complete replacement content..."
      }
    ]
  }
}
\end{lstlisting}

\subsection{HTTP Engine API}

\apitabletitle{Engine HTTP Routes}
\begin{center}
\small
\begin{tabular}{@{}p{0.10\linewidth}p{0.17\linewidth}p{0.22\linewidth}p{0.45\linewidth}@{}}
\toprule
\textbf{Method} & \textbf{Path} & \textbf{Auth} & \textbf{Behavior}\\
\midrule
GET & \apicode{/health} & none & Return \apicode{ok}.\\
POST & \apicode{/hooks} & \apicode{X-Gitlab-Token} matching \apicode{JERYU\_WEBHOOK\_SECRET} & Ingest GitLab job, pipeline, push, and MR webhooks.\\
GET & \apicode{/cache/summary} & none & Return cache health summary JSON.\\
\bottomrule
\end{tabular}
\end{center}

\begin{lstlisting}[style=api,language=bash,caption={HTTP route examples}]
curl http://127.0.0.1:9777/health

curl http://127.0.0.1:9777/cache/summary

curl -X POST http://127.0.0.1:9777/hooks \
  -H "X-Gitlab-Token: $JERYU_WEBHOOK_SECRET" \
  -H "X-Gitlab-Event: Push Hook" \
  --data-binary @push-hook.json
\end{lstlisting}

\apitabletitle{Webhook Event Behavior}
\begin{center}
\small
\begin{tabular}{@{}p{0.20\linewidth}p{0.74\linewidth}@{}}
\toprule
\textbf{Event} & \textbf{Behavior}\\
\midrule
Job Hook & Upsert job event, maybe auto-retry failures, scale up for pending/created work.\\
Pipeline Hook & Track pipeline state, launch canary on green main, update release/prod pipeline state.\\
Push Hook & Normalize ref, skip \apicode{jeryu-test-*}, cancel superseded pipelines, compute impact, persist VTI plan.\\
Merge Request Hook & Logged for observability; no privileged merge behavior.\\
Other & Logged as unhandled.\\
\bottomrule
\end{tabular}
\end{center}

\subsection{Action Registry Contract}

The action registry is the shared action catalog consumed by the TUI command palette, \apicode{jeryu action list}, and \apicode{ListAllowedActions}. Each entry exposes \apicode{id}, \apicode{label}, \apicode{key\_hint}, \apicode{risk\_tier}, \apicode{side\_effect\_class}, \apicode{required\_grant}, \apicode{dry\_run}, \apicode{surfaces}, and \apicode{description}.

\apitabletitle{Canonical Action Registry}
\begin{center}
\small
\begin{tabular}{@{}p{0.22\linewidth}p{0.13\linewidth}p{0.14\linewidth}p{0.16\linewidth}p{0.29\linewidth}@{}}
\toprule
\textbf{Action ID} & \textbf{Risk} & \textbf{Grant} & \textbf{Surfaces} & \textbf{Purpose}\\
\midrule
\apicode{open\_logs} & read-only & none & TUI & Open selected job logs.\\
\apicode{retry\_job} & low & agent task & CLI, TUI & Retry a failed or canceled job.\\
\apicode{delete\_record} & low & agent task & TUI & Remove local DB record.\\
\apicode{pause\_pool} & low & agent task & CLI, TUI & Pause or resume runner pool.\\
\apicode{explain\_blockers} & read-only & none & CLI, TUI, capability & Explain job/release/merge blockers.\\
\apicode{get\_system\_snapshot} & read-only & none & CLI, capability & Return system summary.\\
\apicode{propose\_patch} & high & agent task & CLI, TUI, capability & Create branch, patch, and MR.\\
\apicode{race\_patches} & high & agent task & CLI, capability & Launch competing hypotheses.\\
\apicode{request\_merge} & production & merge approval & CLI, TUI, capability & Request merge through gate.\\
\apicode{plan\_validation} & read-only & none & CLI, capability & Validate proposed test plan.\\
\apicode{run\_tests} & low & agent task & CLI, TUI, capability & Trigger targeted test pipeline.\\
\apicode{next\_action} & read-only & none & CLI, TUI & Print recommended next action.\\
\apicode{tab\_*} & read-only & none & TUI & Switch primary TUI tab from Mission through Secrets.\\
\apicode{toggle\_audit\_ledger} & read-only & none & TUI & Toggle Evidence tab view.\\
\apicode{quit} & read-only & none & TUI & Exit TUI.\\
\bottomrule
\end{tabular}
\end{center}

\begin{lstlisting}[style=api,language=json,caption={Action registry JSON excerpt}]
{
  "id": "request_merge",
  "label": "Request merge",
  "risk_tier": "production",
  "side_effect_class": "merge",
  "required_grant": "merge_approval",
  "dry_run": false,
  "surfaces": ["cli", "tui", "capability"],
  "description": "Merge an MR through the risk gate..."
}
\end{lstlisting}

\clearpage
\twocolumn
\flushbottom
\subsection{Native MCP Server}
\label{app:mcp-server}

V3.01 ships a native Model Context Protocol server for external coding agents that already speak MCP. The server is implemented as \apicode{src/mcp.rs} plus the \apicode{src/mcp/} module tree and is exposed through \apicode{jeryu mcp serve}, \apicode{jeryu mcp serve-http}, and \apicode{jeryu mcp tools --json}. Its purpose is deliberately narrow: translate MCP lifecycle and tool messages into the same typed \apicode{AgentIntent} execution path used by the Unix-socket capability API.

The important boundary is that MCP is a transport, not a privilege escalation path. A call to \apicode{jeryu.request\_merge} over MCP is evaluated by the same capability policy, grant metadata, durable evidence recording, GitLab client wrapper, and merge/release gate logic as the equivalent native capability request. Read-only tools remain read-only; branch-writing tools still create explicit branches and evidence; production-sensitive tools still require the same gate decisions.

\begin{table}[t]
\centering
\tiny
\caption{MCP server transports and lifecycle}
\begin{tabular}{@{}p{0.34\columnwidth}p{0.56\columnwidth}@{}}
\toprule
\textbf{Surface} & \textbf{Behavior}\\
\midrule
\apicode{stdio} & JSON-RPC line protocol for local agents launched as child processes by editors, terminals, or agent runtimes.\\
\apicode{HTTP} & Streamable HTTP endpoint at \apicode{/mcp}, bound to loopback by default through \apicode{settings.mcp.bind}.\\
\apicode{initialize} & Negotiates protocol version \apicode{2025-11-25}, returns \apicode{serverInfo}, and advertises tool support.\\
\shortstack[l]{\apicode{notifications/}\\\apicode{initialized}} & Marks the session initialized and records client attribution for later capability context.\\
\apicode{ping} & Lightweight liveness check for clients that keep long-running sessions.\\
\apicode{tools/list} & Returns the canonical tool manifest derived from the action registry.\\
\apicode{tools/call} & Parses tool arguments, builds an \apicode{AgentIntent}, executes it through capability policy, and returns MCP content plus structured result data.\\
\bottomrule
\end{tabular}
\end{table}

\begin{table}[t]
\centering
\tiny
\caption{MCP tool surface}
\begin{tabular}{@{}p{0.45\columnwidth}p{0.45\columnwidth}@{}}
\toprule
\textbf{Tool} & \textbf{Function}\\
\midrule
\shortstack[l]{\apicode{jeryu.get\_system}\\\apicode{\_snapshot}} & Read current operational state for jobs, pools, release, and evidence.\\
\shortstack[l]{\apicode{jeryu.explain}\\\apicode{\_blockers}} & Explain why a job, release, or merge is blocked.\\
\shortstack[l]{\apicode{jeryu.plan}\\\apicode{\_validation}} & Produce a VTI-backed validation plan for selected tests and refs.\\
\shortstack[l]{\apicode{jeryu.fetch}\\\apicode{\_capsule}} & Fetch structured failure capsule data for a job.\\
\shortstack[l]{\apicode{jeryu.run}\\\apicode{\_tests}} & Create an ephemeral validation branch and trigger CI for a scoped test run.\\
\shortstack[l]{\apicode{jeryu.propose}\\\apicode{\_patch}} & Create a branch, apply file modifications, commit, and open a merge request.\\
\shortstack[l]{\apicode{jeryu.race}\\\apicode{\_patches}} & Launch multiple patch hypotheses while preserving branch-level evidence.\\
\shortstack[l]{\apicode{jeryu.request}\\\apicode{\_merge}} & Ask the merge gate to evaluate an MR against release and risk policy.\\
\bottomrule
\end{tabular}
\end{table}

Tool discovery is schema-bearing. Each MCP tool includes an \apicode{inputSchema}, an \apicode{outputSchema}, and MCP annotations such as \apicode{readOnlyHint}, \apicode{destructiveHint}, \apicode{idempotentHint}, and \apicode{openWorldHint}. These annotations are derived from the same action-risk vocabulary used by the CLI and TUI registry, so an agent does not need a separate policy document to know which tools are read-only, mutating, or production-sensitive.

The HTTP transport adds local-server hardening around the same core. It refuses non-loopback bind addresses at startup, rejects non-local \apicode{Origin} headers, returns \apicode{Mcp-Session-Id} on \apicode{initialize}, requires that session id on later requests, requires \apicode{MCP-Protocol-Version} after initialization, validates \apicode{Mcp-Method} and \apicode{Mcp-Name} against the JSON-RPC body, and supports \apicode{DELETE /mcp} for explicit session teardown. \apicode{GET /mcp} returns \apicode{405}; this implementation is a POST-oriented local control surface rather than an event-streaming API.

\begin{lstlisting}[style=api,language=json,caption={MCP tool call shape}]
{
  "jsonrpc": "2.0",
  "id": 7,
  "method": "tools/call",
  "params": {
    "name": "jeryu.explain_blockers",
    "arguments": {
      "entity_type": "merge",
      "entity_id": 42
    }
  }
}
\end{lstlisting}

The response shape is intentionally useful to both simple chat agents and structured orchestrators. MCP \apicode{content} carries a human-readable text summary, while \apicode{structuredContent} carries the native capability response with \apicode{success}, \apicode{message}, and optional data. When the underlying capability call fails, the server sets \apicode{isError} without hiding the structured policy result. This lets clients present a concise explanation to a human while still preserving machine-readable evidence for retries, escalation, or audit.
